\section{Steering Dynamics \& Self-Aligning Torque}

\label{sec:steering-sat}

This section documents the steering column dynamics and self-aligning torque (SAT), which make the steering wheel feel responsive and natural.

\cite{pacejka2005tire} provides the theoretical foundation for self-aligning torque. The implementation uses a simplified column model with spring return and viscous damping.

\subsection{Self-Aligning Torque (SAT)}

Self-aligning torque is the tendency of tires to straighten when cornering, caused by asymmetric pressure distribution:

\begin{equation}
  \text{SAT} = -k_r \cdot d \cdot F_{\text{lat}}
  \label{eq:sat-magnitude}
\end{equation}

where:
\begin{itemize}
  \item $k_r$ = pneumatic trail (moment arm, typically 0.03–0.05 m)
  \item $d$ = lateral force direction (sign of sideways slip)
  \item $F_{\text{lat}}$ = lateral tire force (magnitude)
\end{itemize}

A simplified version sums SAT from front wheels only (dominant contribution):

\begin{equation}
  \text{SAT}_{\text{total}} = -k_r \cdot (F_{\text{lat,FL}} + F_{\text{lat,FR}}) \cdot \frac{1}{w}
  \label{eq:sat-total}
\end{equation}

where $w$ is the front track width.

\textbf{Source Code}: \coderef{steering.js}{All}

\textbf{Notes}:
\begin{itemize}
  \item SAT increases with cornering force (stronger feedback at higher speeds)
  \item SAT decreases at high slip angles (sliding tires have less effect)
  \item Only front wheels contribute significantly
  \item The pneumatic trail $k_r$ is a tire property, tuned for responsiveness
\end{itemize}

\subsection{Steering Column Equation of Motion}

The steering column dynamics are modeled as:

\begin{equation}
  I_c \cdot \ddot{\delta} = \text{SAT} - c_v \cdot \dot{\delta} - c_c \cdot \text{sign}(\dot{\delta}) + F_s \cdot (\delta_{\text{target}} - \delta)
  \label{eq:steering-column}
\end{equation}

where:
\begin{itemize}
  \item $I_c$ = steering column inertia (kg·m²)
  \item $\delta$ = steering angle (radians)
  \item $c_v$ = viscous damping coefficient
  \item $c_c$ = Coulomb (friction) damping
  \item $F_s$ = centering spring stiffness
  \item $\delta_{\text{target}}$ = driver's requested steering angle
\end{itemize}

\subsection{Semi-Implicit Euler Integration}

For numerical stability, the steering column is integrated using semi-implicit Euler:

\begin{equation}
  \omega_{\text{new}} = \frac{\omega_{\text{old}} + (\text{SAT} - c_v \omega_{\text{old}}) / I_c \cdot \Delta t}{1 + c_v \Delta t / I_c}
  \label{eq:steering-semi-implicit}
\end{equation}

\begin{equation}
  \delta_{\text{new}} = \delta_{\text{old}} + \omega_{\text{new}} \cdot \Delta t
  \label{eq:steering-angle-update}
\end{equation}

\textbf{Physical Meaning}: The semi-implicit formulation provides better stability at high SAT values without requiring a smaller timestep.

\textbf{Source Code}: \coderef{steering.js}{Implementation details}

\subsection{Low-Speed Spring Return}

At low speeds (below 3 m/s), SAT is reduced and a spring return force assists steering:

\begin{equation}
  F_{\text{return}} = -k_s \cdot (1 - \text{speedFactor}) \cdot \delta
  \label{eq:spring-return}
\end{equation}

where:

\begin{equation}
  \text{speedFactor} = \max(0, \min(1, v / 3.0))
  \label{eq:speed-factor}
\end{equation}

\textbf{Physical Meaning}: At standstill, the spring provides full return force. As speed increases, the spring force decreases and SAT takes over. This ensures the wheel self-centers when parked.

\textbf{Notes}:
\begin{itemize}
  \item Spring stiffness $k_s$ is tuned for feel; typical value 100–200 N/m
  \item Transition range: 0–3 m/s
  \item SAT feedback strengthens above 3 m/s
\end{itemize}

\newpage
